%!TEX root = ../memoria.tex
% =============================================================================
%  APÉNDICE A — Detalle de la evaluación de usabilidad (AIPO)
%
%  Este apéndice complementa la Sección 4.4 con el material detallado
%  de las dos sesiones de evaluación. Se incluye:
%    A.1 — Protocolo de las tareas (instrucciones literales).
%    A.2 — Tablas de incidencias por tarea (P1 y P2).
%    A.3 — Cuestionario final y valoraciones.
% =============================================================================

Este apéndice recoge el material completo de las dos sesiones de
evaluación con personas mayores descritas en la
Sección~\ref{SEC:IMP_AIPO}. Se han trasladado a esta ubicación
los elementos que, sin formar parte del análisis principal,
permiten verificar el procedimiento seguido y reproducir la
evaluación en futuras iteraciones del proyecto.


\section{Protocolo de las
tareas\label{SEC:APP_PROTOCOLO}}

Cada sesión ha comenzado con una breve introducción al
participante, en la que se ha explicado el carácter de la prueba,
se ha solicitado el consentimiento verbal y se ha indicado al
participante que podía verbalizar en cualquier momento sus dudas
o impresiones. A continuación se han planteado las cinco tareas
en el orden que se recoge a continuación. La redacción literal de
las instrucciones se ha mantenido lo más natural posible para no
introducir vocabulario técnico que pudiera condicionar el
resultado.

\subsection{Tarea T1 — Añadir un recordatorio}

\emph{``Imagine que tiene que ir al médico el viernes que viene a
las once de la mañana. Apunte esa cita en la aplicación para que
no se le olvide''.}

Antes de comenzar la tarea, se ha preguntado al participante
cómo crearía su propio perfil de usuario en la aplicación, con
el fin de explorar la comprensión de los conceptos de
\emph{nombre de usuario} y \emph{código de grupo familiar}.

\subsection{Tarea T2 — Marcar y modificar un recordatorio}

\emph{``Acaba de tomar la pastilla de la tensión que aparece en
la pantalla principal. Indique en la aplicación que ya la ha
tomado''.}

Como ampliación, se ha preguntado al participante cómo deshacer
un recordatorio marcado por error y cómo modificar el título o
la hora de uno ya existente.

\subsection{Tarea T3 — Llamar a un contacto y dar de alta uno
nuevo}

\emph{``Su hija le ha pedido que llame al médico de cabecera.
Encuentre su teléfono en la aplicación e inicie la llamada''.}

Como ampliación, se ha pedido al participante que añadiera un
contacto nuevo —por ejemplo, el teléfono de un vecino— para
evaluar la comprensión del formulario de contactos y de las
categorías disponibles.

\subsection{Tarea T4 — Enviar un mensaje}

\emph{``Quiere avisar a su hija de que ya ha ido al médico.
Envíele un mensaje desde la aplicación y lea la respuesta cuando
le llegue''.}

\subsection{Tarea T5 — Consultar la galería y subir una
fotografía}

\emph{``Su hija le ha mandado unas fotos del nieto. Ábralas y
déjele un comentario en una de ellas''.}

Como ampliación, se ha pedido al participante que subiera una
fotografía propia desde la \emph{tablet} para comprobar la
comprensión del formulario de subida y del selector de archivos.


\section{Tablas de incidencias por
tarea\label{SEC:APP_INCIDENCIAS}}

Para cada tarea se ha registrado, durante la propia sesión, la
información que se recoge en las tablas siguientes: la duración
aproximada de la tarea, su resultado —completada, completada con
ayuda o no completada— y las incidencias observadas, con su
gravedad estimada.

La gravedad se clasifica en tres niveles, conforme a la
recomendación de Nielsen para evaluaciones heurísticas: 
\textbf{baja} —el problema causa molestia menor pero no impide
completar la tarea—, \textbf{media} —el problema requiere
exploración adicional o produce confusión— y \textbf{alta} —el
problema impide completar la tarea sin ayuda externa—.

\begin{table}[incidencias T1]{TAB:APP_T1}{Resultados e
incidencias de la tarea T1 (añadir un recordatorio).}
\begin{center}
\begin{tabular}{p{3.5cm}p{4.5cm}p{4.5cm}}
\toprule
& \textbf{P1} & \textbf{P2} \\
\midrule
Duración aproximada    & 3 minutos      & 2 -- 4 minutos \\
Resultado              & Completada con dificultad & Completada con dificultad \\
\midrule
Incidencias  &
\begin{itemize}\setlength\itemsep{0pt}\setlength\parsep{0pt}
\item Confusión sobre el concepto de \emph{nombre de usuario} (alta).
\item El campo de fecha no aparecía al principio del formulario (media).
\item El campo de hora figuraba como obligatorio (media).
\item Botón \emph{Añadir recordatorio} poco visible (baja).
\end{itemize}
&
\begin{itemize}\setlength\itemsep{0pt}\setlength\parsep{0pt}
\item Dificultad para localizar el botón de creación (media).
\item Confusión sobre el orden de los campos del formulario (media).
\item Duda sobre si era obligatorio rellenar la hora (baja).
\end{itemize}\\
\bottomrule
\end{tabular}
\end{center}
\end{table}
\FloatBarrier

\begin{table}[incidencias T2]{TAB:APP_T2}{Resultados e
incidencias de la tarea T2 (marcar y modificar un recordatorio).}
\begin{center}
\begin{tabular}{p{3.5cm}p{4.5cm}p{4.5cm}}
\toprule
& \textbf{P1} & \textbf{P2} \\
\midrule
Duración aproximada    & 1 minuto       & Menos de 2 minutos \\
Resultado              & Completada     & Completada \\
\midrule
Incidencias &
\begin{itemize}\setlength\itemsep{0pt}\setlength\parsep{0pt}
\item Confusión entre el botón de posponer y el de editar (media).
\item Botón \emph{Deshacer} poco visible (baja).
\item Recordatorios pendientes y completados no se diferenciaban visualmente (baja).
\end{itemize}
&
\begin{itemize}\setlength\itemsep{0pt}\setlength\parsep{0pt}
\item Confusión entre el botón de posponer y el de editar (media).
\item Identificó sin problemas el botón \emph{Hecho} (sin incidencia).
\end{itemize}\\
\bottomrule
\end{tabular}
\end{center}
\end{table}
\FloatBarrier

\begin{table}[incidencias T3]{TAB:APP_T3}{Resultados e
incidencias de la tarea T3 (llamar y dar de alta un contacto).}
\begin{center}
\begin{tabular}{p{3.5cm}p{4.5cm}p{4.5cm}}
\toprule
& \textbf{P1} & \textbf{P2} \\
\midrule
Duración aproximada    & 1 -- 2 minutos & 1 -- 2 minutos \\
Resultado              & Completada     & Completada \\
\midrule
Incidencias &
\begin{itemize}\setlength\itemsep{0pt}\setlength\parsep{0pt}
\item Intentó utilizar el botón \emph{atrás} del navegador, sin éxito (media).
\item Botón de edición del contacto poco evidente (baja).
\item Sin validación del formato del teléfono al dar de alta (baja).
\end{itemize}
&
\begin{itemize}\setlength\itemsep{0pt}\setlength\parsep{0pt}
\item Localizó el botón \emph{Llamar} sin ayuda (sin incidencia).
\item Duda sobre en qué categoría clasificar el contacto nuevo (baja).
\end{itemize}\\
\bottomrule
\end{tabular}
\end{center}
\end{table}
\FloatBarrier

\begin{table}[incidencias T4]{TAB:APP_T4}{Resultados e
incidencias de la tarea T4 (enviar un mensaje).}
\begin{center}
\begin{tabular}{p{3.5cm}p{4.5cm}p{4.5cm}}
\toprule
& \textbf{P1} & \textbf{P2} \\
\midrule
Duración aproximada    & 2 -- 4 minutos & 2 -- 4 minutos \\
Resultado              & Completada con ayuda & Completada con ayuda \\
\midrule
Incidencias &
\begin{itemize}\setlength\itemsep{0pt}\setlength\parsep{0pt}
\item Campo de escritura poco visible, en la parte inferior de la pantalla (alta).
\item En el tema oscuro, las burbujas pierden contraste (media).
\item Las letras del chat resultan algo pequeñas (baja).
\end{itemize}
&
\begin{itemize}\setlength\itemsep{0pt}\setlength\parsep{0pt}
\item Campo de escritura poco visible (alta).
\item Identificó las burbujas propias y recibidas sin ayuda (sin incidencia).
\end{itemize}\\
\bottomrule
\end{tabular}
\end{center}
\end{table}
\FloatBarrier

\begin{table}[incidencias T5]{TAB:APP_T5}{Resultados e
incidencias de la tarea T5 (consultar galería y subir una foto).}
\begin{center}
\begin{tabular}{p{3.5cm}p{4.5cm}p{4.5cm}}
\toprule
& \textbf{P1} & \textbf{P2} \\
\midrule
Duración aproximada    & 2 -- 4 minutos & 2 -- 4 minutos \\
Resultado              & Completada con ayuda & Completada con ayuda \\
\midrule
Incidencias &
\begin{itemize}\setlength\itemsep{0pt}\setlength\parsep{0pt}
\item No identificó cómo ampliar una foto a pantalla completa (media).
\item Texto del selector de archivos interpretado como error (alta).
\item Identificó sin problemas el botón \emph{Subir foto} (sin incidencia).
\end{itemize}
&
\begin{itemize}\setlength\itemsep{0pt}\setlength\parsep{0pt}
\item Mismo problema con el texto del selector de archivos (alta).
\item Duda sobre si una foto subida era visible para el resto del grupo (baja).
\end{itemize}\\
\bottomrule
\end{tabular}
\end{center}
\end{table}
\FloatBarrier


\section{Cuestionario final y
valoraciones\label{SEC:APP_CUESTIONARIO}}

Al finalizar las cinco tareas, se ha planteado a cada participante
un breve cuestionario que combina valoraciones cuantitativas
—mediante escala Likert de uno a cinco— con preguntas abiertas
sobre la experiencia general. Los resultados se recogen en la
Tabla~\ref{TAB:APP_CUESTIONARIO}.

\begin{table}[cuestionario final]{TAB:APP_CUESTIONARIO}{Valoraciones
del cuestionario final por participante.}
\begin{center}
\begin{tabular}{p{6cm}cc}
\toprule
\textbf{Aspecto valorado} & \textbf{P1} & \textbf{P2} \\
\midrule
Facilidad de uso general (1--5)       & 5  & 4 \\
Comprensión de los iconos (1--5)      & 4  & 4 \\
Tamaño de los elementos (1--5)        & 4  & 4 \\
Utilidad percibida (1--5)             & 4  & 5 \\
Confianza al utilizar la aplicación (1--5) & 5 & 4 \\
\midrule
\multicolumn{3}{l}{\emph{Preguntas abiertas}} \\
\midrule
Módulo preferido & Mensajes, contactos y fotos & Contactos \\
Módulo más difícil & Agenda & Agenda \\
¿Utilizaría la aplicación? & Sí & Sí \\
\bottomrule
\end{tabular}
\end{center}
\end{table}
\FloatBarrier

Las valoraciones son consistentes entre ambos participantes y
coherentes con las incidencias observadas durante las tareas: la
mayor dificultad se concentra en la agenda —donde los formularios
de creación son los más complejos del sistema— mientras que los
módulos basados en la consulta o en interacciones directas
—contactos, mensajes y galería— han recibido valoraciones más
altas. Ambos participantes han manifestado interés en seguir
utilizando la aplicación más allá de la sesión de prueba, lo que
constituye una validación cualitativa de la utilidad práctica de
\emph{ConectaMayor}.
